%% LyX 2.4.4 created this file.  For more info, see https://www.lyx.org/.
%% Do not edit unless you really know what you are doing.
\documentclass[english,t]{beamer}
\usepackage{mathptmx}
\usepackage[T1]{fontenc}
\usepackage[latin9]{inputenc}
\usepackage{amssymb}
\usepackage{cancel}

\makeatletter

%%%%%%%%%%%%%%%%%%%%%%%%%%%%%% LyX specific LaTeX commands.
%% Because html converters don't know tabularnewline
\providecommand{\tabularnewline}{\\}

%%%%%%%%%%%%%%%%%%%%%%%%%%%%%% Textclass specific LaTeX commands.
% this default might be overridden by plain title style
\newcommand\makebeamertitle{\frame{\maketitle}}%
% (ERT) argument for the TOC
\AtBeginDocument{%
  \let\origtableofcontents=\tableofcontents
  \def\tableofcontents{\@ifnextchar[{\origtableofcontents}{\gobbletableofcontents}}
  \def\gobbletableofcontents#1{\origtableofcontents}
}

%%%%%%%%%%%%%%%%%%%%%%%%%%%%%% User specified LaTeX commands.
\usetheme{Warsaw}
% or ...

\setbeamercovered{transparent}
% or whatever (possibly just delete it)
\setbeamercovered{invisible} %makes overlays invisible

%gets rid of navigation symbols
\setbeamertemplate{navigation symbols}{}

%\usepackage{hyperref}
%\usepackage[usenames,dvipsnames,table]{xcolor} %adds additional color options
\definecolor{links}{HTML}{2A1B81}
%\hypersetup{colorlinks,linkcolor=blue,urlcolor=links}

\usepackage{multicol}

\usepackage{hyperref}

\usepackage{tikz}
\usetikzlibrary{calc}
\usetikzlibrary{plotmarks}
\usetikzlibrary{shapes}
\usepackage{pgfplots}
\pgfplotsset{compat=newest}

\usepackage{forest} %%for tree diagram

\usepackage{calc}
\usepackage[nomessages]{fp}

\usepackage[super]{nth} %easy 1st, 2nd, etc

%\usepackage{advdate}    % Advancing/saving dates
\usepackage{datetime}   % Dates formatting
%\usepackage{datenumber} % Counters for dates

\usepackage[super]{nth}

\usepackage{etex}

\usepackage{fancybox}

%%table colors
%\usepackage[table]{xcolor}
\usepackage{colortbl}

\usepackage[outline]{contour}  %allows text halos

\usepackage{epsdice} %%dice font
\usepackage[misc]{ifsym}
\newfont\dice{dice3d} %%3d dice font

\contourlength{1pt} %sets halo width
%%example: \node at (0.75,0.5) {\contour{white}{\Large Text!}};
%%example:  \node [fill=white,inner sep=1pt] at (2.25,0.5) {\Large 			Text!};
%%example:  \node [fill=white,rounded corners=2pt,inner sep=1pt] at 			(3.75,0.5) {\Large Text!};

%%%%%%%%%%%%%%%%%%%%%%%%%%%
\def\formatdateny{\csname noyear\languagename\expandafter\endcsname\formatdate}
\def\noyearenglish#1, \the\@year{#1}
\let\noyearamerican\noyearenglish
\let\noyearbritish\noyearenglish

%%allows uncover with forest diagrams
\tikzset{
    invisible/.style={opacity=0,text opacity=0},
    visible on/.style={alt=#1{}{invisible}},
    alt/.code args={<#1>#2#3}{%
      \alt<#1>{\pgfkeysalso{#2}}{\pgfkeysalso{#3}} % \pgfkeysalso doesn't change the path
    },
}
\forestset{
  visible on/.style={
    for tree={
      /tikz/visible on={#1},
      edge={/tikz/visible on={#1}}}}}

%%%redefines column alignment to match vertically
%\define@key{beamerframe}{t}[true]{% top
  %\beamer@frametopskip=.2cm plus .5\paperheight\relax%
  %\beamer@framebottomskip=0pt plus 1fill\relax%
 % \beamer@frametopskipautobreak=\beamer@frametopskip\relax%
  %\beamer@framebottomskipautobreak=\beamer@framebottomskip\relax%
  %\def\beamer@initfirstlineunskip{}%
%}

%%provides pdf with 4 slides per page; don't forget to add 'handout' to remove overlays
%\usepackage{pgfpages}
%\pgfpagesuselayout{4 on 1}[a4paper, landscape, border shrink=7mm]
%\pgfpageslogicalpageoptions{1}{border code=\pgfusepath{stroke}}
%\pgfpageslogicalpageoptions{2}{border code=\pgfusepath{stroke}}
%\pgfpageslogicalpageoptions{3}{border code=\pgfusepath{stroke}}
%\pgfpageslogicalpageoptions{4}{border code=\pgfusepath{stroke}}
%%%Don't forget to add 'handout'

\makeatother

\usepackage{babel}
\begin{document}
\newcommand{\xmax}{2000}
\newcommand{\xmin}{0}
\newcommand{\xstep}{0} %must be xmin+n
\newcommand{\ymax}{500}
\newcommand{\ymin}{0}
\newcommand{\ystep}{0} %must be ymin+n

\newcounter{countEx} \setcounter{countEx}{0}
\newcounter{countProb} \setcounter{countProb}{0}
\resetcounteronoverlays{countEx} \resetcounteronoverlays{countProb}

\newcommand{\ex}[3]{
	\addtocounter{countEx}{1}
	\only<#1-#2>{\begin{exampleblock} {Example \chNum.\secNum.\arabic{countEx}.} #3	\end{exampleblock}}
}

\newcommand{\prob}[4]{
	\addtocounter{countProb}{1}
	\only<#1-#2>{\begin{exampleblock} {Problem  \chNum.\secNum.\arabic{countProb}.\,#3} #4	\end{exampleblock}}
}

\newcommand{\yline}[4]{
	(#2 - #4)/(#1 - #3)*(x - #1) + #2
}

\beamerdefaultoverlayspecification{<+->}

%%%%Chapter and Section numbers%%%%%
\newcommand{\chNum}{3}
\newcommand{\secNum}{5}
\newcommand{\secTitle}{Expected Value}
\date{\formatdate{17}{10}{2014}}

\title[Section \chNum.\secNum: \secTitle]{Section \chNum.\secNum: \secTitle}

\subtitle{Finite Mathematics~~MTH 107~~Fall 2014}

\author{Dr.\,James Ashe}

\titlegraphic{\includegraphics[height=2.75cm]{/home/jim/JamesAshe/MHU/images/MarsHilUniversityLogo-Virtical.png}}

\institute{Mars Hill University}

\makebeamertitle 

%%True false command
\newcommand{\T}{\textcolor{green}{\contour{black}{T}}}
\newcommand{\F}{\textcolor{red}{\contour{black}{F}}}
\newcommand{\uT}[2]{\uncover<#1-#2>{\T}}
\newcommand{\uF}[2]{\uncover<#1-#2>{\F}}

%tkiz ball item 
\newcommand*\circled[1]{\tikz[baseline=(char.base)]
	{\node[circle,ball color=green!50!black!100, shade,   color=white,inner sep=1.2pt] (char) {\tiny #1};
	}
}
\begin{frame}{}

\vspace{-.5cm}
\begin{exampleblock}{Average Value}
John has homework grades of 70\%, 80\%, 50\%, and 100\%. What is
his \emph{average} grade?

\begin{itemize}
\item <2->[]$\dfrac{.70+.80+.50+.100}{4}=75\%$
\end{itemize}
\end{exampleblock}
\begin{itemize}
\item <3->A summary of what \emph{has }happened. 
\end{itemize}
\begin{exampleblock}<4->{Expected Value}
John places bets of \$70, \$80, \$50, and \$100 on a game in which
he has a 25\% chance of winning. How much money should he expect to
win? 

\begin{itemize}
\item <5->[]$70\cdot\dfrac{1}{4}+80\cdot\dfrac{1}{4}+50\cdot\dfrac{1}{4}+100\cdot\dfrac{1}{4}=\$75$
\end{itemize}
\end{exampleblock}
\begin{itemize}
\item <6->A summary of what \emph{will }happen.
\end{itemize}
\end{frame}

\begin{frame}{Expected Value}

\vspace{-.5cm}
\begin{block}{Expected Value (``long-term average'') To find the \textbf{expected
value }of an experiment, multiply the \emph{value} of each outcome
by its probability and add the results.}
\end{block}
\begin{columns}


\column{8cm}

\ex{2}{}{A saleswoman has found that her weekly commissions have
the following probabilities.
\begin{itemize}
\item []\scriptsize%
\begin{tabular}{|c|c|c|c|c|}
\hline 
\cellcolor{yellow!50}\textbf{commision} & $0$ & $\$100$ & $\$200$ & $\$300$\tabularnewline
\hline 
\cellcolor{yellow!50}\textbf{probability} & $0.05$ & $0.35$ & $0.50$ & $0.10$\tabularnewline
\hline 
\end{tabular}\normalsize
\end{itemize}
}
\end{columns}
\begin{itemize}
\item []<3->\textbf{Expected Commission:} \uncover<4->{$=(0)(.05)$}\uncover<5->{$+(100)(.35)$}\uncover<6->{$+(200)(.50)$}\uncover<7->{$+(300)(.10)$}\uncover<8->{$=\$165$}
\item <9->Based on her sales history, she should expect to make $\$165$
a week in commissions. 
\end{itemize}
\end{frame}

\begin{frame}{Decisions}

\ex{1}{}{You are asked to play the following game:
\begin{itemize}
\item If you roll a 1 or 2, you win $\$50$.
\item If you roll a 3, you lose $\$20$.
\item If you roll a 4, 5, or 6, you lose $\$30$.
\end{itemize}
Should you play?}
\begin{columns}


\column{5cm}
\begin{itemize}
\item []<2->$p(1\cup2)=$\uncover<3->{ $\frac{1}{6}+\frac{1}{6}=\frac{1}{3}$}
\item []<3->$p(3)=$\uncover<4->{ $\frac{1}{6}$}
\item []<5->$p(4\cup5\cup6)$\uncover<6->{ \hspace{0cm}\hphantom{HI}$=\frac{1}{6}+\frac{1}{6}+\frac{1}{6}=\frac{1}{2}$}
\end{itemize}

\column{6cm}
\begin{itemize}
\item []<7->The expected value of the game, 
\item []<8->$EV=$\uncover<9->{ $50\cdot\frac{1}{3}$}\uncover<10->{ $-20\cdot\frac{1}{6}$}\uncover<11->{ $-30\cdot\frac{1}{2}$}
\item []<12->\hspace{0cm}\hphantom{$EV$}\uncover<13->{ $\approx-\$1.67$}
\item []<14>So NO.
\end{itemize}
\end{columns}
\end{frame}

\begin{frame}{}

\begin{columns}


\column{6cm}

\ex{1}{}{A club is having a raffle. Tickets cost $\$5$, but each
prize will be given away regardless of tickets sold. Should you buy
a ticket under the following circumstances:
\begin{itemize}
\item [a.]<2->1000 tickets are sold. 
\item [b.]<11->700 tickets are sold.
\end{itemize}
}

\column{5.5cm}

\vspace{1cm}\normalsize%
\begin{tabular}{|c|c|c|}
\hline 
\rowcolor{yellow!50}\textbf{Prize} & \textbf{Value} & \textbf{\# }\tabularnewline
\hline 
Laptop & $\$2,325$ & 1\tabularnewline
\hline 
MP3 Player & $\$425$ & 2\tabularnewline
\hline 
20 CDs & $\$320$ & 3\tabularnewline
\hline 
A Pizza & $\$23$ & 4\tabularnewline
\hline 
Book & $\$20$ & \hphantom{$5$}15\tabularnewline
\hline 
\multicolumn{1}{c}{\uncover<7->{nothing}} & \multicolumn{1}{c}{\uncover<7->{$\$0$}} & \multicolumn{1}{c}{\only<9-10>{$975$}\only<13->{$675$}}\tabularnewline
\end{tabular}\normalsize
\end{columns}
\begin{itemize}
\item <3->Subtract $\$5$ from the value of the prizes for the ticket price.
\item <8->There are \only<8-10>{$1000-25=975$ losing tickets.}\only<13->{$700-25=675$
losing tickets.}
\item <3->[a.]\scriptsize $EV=$\uncover<4->{ $2320\left(\frac{1}{1000}\right)$}\uncover<5->{$+420\left(\frac{2}{1000}\right)+315\left(\frac{3}{1000}\right)+18\left(\frac{4}{1000}\right)+15\left(\frac{15}{1000}\right)$}\uncover<10->{$-5$$\left(\frac{975}{1000}\right)$}\uncover<10->{ \hspace{0cm}\hphantom{$EV$} $\approx-\$0.473$}\vspace{-.25cm}
\end{itemize}

\begin{itemize}
\item <12->[b.]\scriptsize $EV=$\uncover<12->{ $2320\left(\frac{1}{700}\right)$}\uncover<12->{ $+420\left(\frac{2}{700}\right)+315\left(\frac{3}{700}\right)+18\left(\frac{4}{700}\right)+15\left(\frac{15}{700}\right)$}\uncover<14->{$-5$$\left(\frac{675}{1000}\right)\approx\$1.10$}
\end{itemize}
\end{frame}

\end{document}
